\section{Deep-Unfolded Receiver Design}
\label{sec:algorithm}

We develop a deep-unfolded MP receiver with a score-based channel
prior for joint activity detection and channel estimation. The
likelihood and priors in \eqref{eq:working-posterior} provide a
starting point for constructing its estimation modules. Following
\cite{cai2026}, a linear estimator processes the received pilots,
while a channel and activity denoiser incorporates the prior
information; the two modules exchange Gaussian extrinsic messages. To adapt the denoisers to the inputs
produced by this receiver, we unfold the message updates into $L$
layers with trainable channel denoisers, activity classifiers, and
message parameters. The linear estimator uses the sensing matrix
formed from the pilots and inter-frame phase increments. The nonlinear
module estimates each user's conditional channel and activity
probability, whose product determines the feedback to the next layer.
Section~\ref{sec:learning} trains these coupled updates using received
pilot observations. In the following,
$J=KN$ and $C=MN$.

\subsection{Linear Estimation}

At layer $\ell$, the input Gaussian messages have means collected in
$\mathbf M_\ell$ and a common covariance $w_\ell\I_J$ for each
column of $\mathbf X$. Combining them with the likelihood in
\eqref{eq:working-posterior} and completing the square gives the
mean $\mathbf M_\ell+\bm\Delta_\ell$ and average marginal variance
$t_\ell$, where
\begin{equation}
 \begin{aligned}
 \bm\Delta_\ell&=w_\ell\mathbf A_{\rm c}^{\mathsf H}
 (\sigma_{\rm c}^2\I_Q+w_\ell\mathbf A_{\rm c}\mathbf A_{\rm c}^{\mathsf H})^{-1}
 (\mathbf Y-\mathbf A_{\rm c}\mathbf M_\ell),
 \\
 t_\ell&=\frac1J\tr\!\left[
 \left(w_\ell^{-1}\I_J+
 \frac{\mathbf A_{\rm c}^{\mathsf H}\mathbf A_{\rm c}}{\sigma_{\rm c}^2}
 \right)^{-1}\right].
 \end{aligned}
 \label{eq:twc-linear-correction}
\end{equation}
Approximating the covariance by $t_\ell\I_J$, we obtain the
extrinsic message by dividing the resulting Gaussian by the input
message \cite{minka2001}. Its variance and mean are given by
\cite{cai2026}
\begin{equation}
 \begin{aligned}
 v_\ell^{(0)}&=(t_\ell^{-1}-w_\ell^{-1})^{-1},\\
 \mathbf R_\ell^{(0)}&=v_\ell^{(0)}
 \left(\frac{\mathbf M_\ell+\bm\Delta_\ell}{t_\ell}
       -\frac{\mathbf M_\ell}{w_\ell}\right).
 \end{aligned}
 \label{eq:twc-linear-extrinsic}
\end{equation}
To adapt these messages during receiver training, we introduce a
positive factor $\gamma_\ell=e^{c_{\gamma,\ell}}$ that scales the
extrinsic variance and adjust the mean as
\begin{equation}
 v_\ell=\gamma_\ell v_\ell^{(0)},\qquad
 \mathbf R_\ell=\mathbf M_\ell+
 \left(1+\frac{v_\ell}{w_\ell}\right)\bm\Delta_\ell.
 \label{eq:twc-query-calibrated}
\end{equation}
For $\gamma_\ell=1$, \eqref{eq:twc-query-calibrated} reduces to
\eqref{eq:twc-linear-extrinsic}. The pair
$(\mathbf R_\ell,v_\ell)$ is then passed to the channel and activity
denoiser.

\subsection{Channel and Activity Denoising}

The Gaussian messages define a local estimation model for constructing
the channel denoiser and activity classifier. For a user block
$r\in\mathbb C^D$ of $\mathbf R_\ell$, the input message is
proportional to $\CN(r;\alpha h,v\I_D)$.
Combining this message with the channel prior in
\eqref{eq:working-prior}, the densities of $r$ under the inactive
and active hypotheses are, respectively,
\begin{equation}
 \begin{aligned}
 q_0(r;v)&=\CN(r;0,v\I_D),\\
 q_1(r;v)&=\int\CN(r;h,v\I_D)\,q_{\rm tr}(dh).
 \end{aligned}
 \label{eq:twc-local-densities}
\end{equation}
Under this local model, the joint channel and activity posterior is
\begin{equation}
 \begin{gathered}
 q(dh,a\mid r,v)=
 \frac{\rho^a(1-\rho)^{1-a}\CN(r;ah,v\I_D)\,q_{\rm tr}(dh)}
 {(1-\rho)q_0(r;v)+\rho q_1(r;v)},\\
 a\in\{0,1\}.
 \end{gathered}
 \label{eq:twc-local-posterior}
\end{equation}
Conditioning \eqref{eq:twc-local-posterior} on the active hypothesis
gives the conditional channel mean
\begin{equation}
 \mu^\star(r,v)=\E[h\mid r,\alpha=1]
 =\frac{\int h\,\CN(r;h,v\I_D)\,q_{\rm tr}(dh)}{q_1(r;v)}.
 \label{eq:twc-local-channel-target}
\end{equation}
The activity probability follows from Bayes' rule as
\begin{equation}
 p^\star(r,v)=\Pr(\alpha=1\mid r)
 =\frac{\rho q_1(r;v)}
 {(1-\rho)q_0(r;v)+\rho q_1(r;v)}.
 \label{eq:twc-local-activity-target}
\end{equation}
Since $x=\alpha h$ vanishes under the inactive hypothesis, averaging
over the two activity states gives \cite{cai2026}
\begin{equation}
 \E[x\mid r]=p^\star(r,v)\mu^\star(r,v).
 \label{eq:twc-mixture-mean}
\end{equation}
Thus channel denoising and activity detection provide the two
quantities needed to estimate the effective channel. Computing them
directly requires integration over the $D$-dimensional complex channel
distribution in \eqref{eq:twc-local-densities} and
\eqref{eq:twc-local-channel-target}. This decomposition motivates a
score-based channel denoiser and an activity classifier whose outputs
are multiplied to form the effective-channel estimate. Their
parameters are then jointly trained through the unfolded iterations
using channel and activity supervision.

Let $u_{k\ell}\in\mathbb R^{2D}$ denote the real and imaginary parts
of the denoiser input for user $k$. Based on the score representation developed
in Section~\ref{sec:learning}, the channel estimate is
\begin{equation}
 \mu_{k\ell}=u_{k\ell}
       +a_\ell\frac{v_\ell}{2}s_{\theta_\ell}(u_{k\ell},v_\ell),
 \qquad a_\ell=e^{c_{\alpha,\ell}}.
 \label{eq:twc-conditional-mean}
\end{equation}
Here $s_{\theta_\ell}$ is the score network in layer $\ell$,
$v_\ell/2$ is the noise variance per real coordinate, and $a_\ell$
controls the denoising correction.

For the classifier, rewriting \eqref{eq:twc-local-activity-target}
in terms of posterior odds gives
\begin{equation}
 \logit p^\star(r,v)
 =\logit(\rho)+\log\frac{q_1(r;v)}{q_0(r;v)}.
 \label{eq:twc-activity-odds}
\end{equation}
The log density ratio supplies the observation-dependent correction
to the prior log odds under the local model. Motivated by this form,
we parameterize the correction with a trainable classifier
\cite{menon2016}.
Using the denoiser input and channel estimate, we compute
\begin{equation}
 p_{k\ell}=\sigmoid\!\left(\logit(\rho)+
 g_{\phi_\ell}(u_{k\ell},v_\ell,\mu_{k\ell},\rho)\right),
 \label{eq:twc-activity}
\end{equation}
where $g_{\phi_\ell}$ parameterizes this correction. Substituting
the channel and activity estimates into
\eqref{eq:twc-mixture-mean} gives
\begin{equation}
 \widehat x_{k\ell}=p_{k\ell}\mu_{k\ell}.
 \label{eq:twc-estimate}
\end{equation}
The user estimates are converted back to complex form and arranged
as $\widehat{\mathbf X}_\ell$ according to \eqref{eq:mmv-indexing}.
Within each layer, all users share the score network and activity
classifier.

\subsection{Extrinsic Message Update}

To form the message returned to the linear estimator, we combine
the mean $\widehat{\mathbf X}_\ell$ with a shared variance estimated
from the observation residual. This supplies the feedback variance
without a second-order score network. We subtract the noise
contribution from the squared residual and normalize it as
\begin{equation}
 e_\ell=\frac{\|\mathbf Y-\mathbf A_{\rm c}\widehat{\mathbf X}_\ell\|_F^2
                   -QC\sigma_{\rm c}^2}
                  {C\|\mathbf A_{\rm c}\|_F^2}.
 \label{eq:twc-residual-scale}
\end{equation}
We then compute the ratio $q_\ell$ of the output variance to the input
variance through
\begin{equation}
 \begin{aligned}
 z_\ell&=e^{c_{\beta,\ell}}
 \left[\eta v_\ell\operatorname{softplus}
       \!\left(\frac{e_\ell}{\eta v_\ell}\right)+\epsilon_s\right],\\
 q_\ell&=\epsilon_r+(q_{\max}-\epsilon_r)
 \tanh\!\left(\frac{z_\ell}{(q_{\max}-\epsilon_r)v_\ell}\right),
 \end{aligned}
 \label{eq:twc-feedback-ratio}
\end{equation}
where $\operatorname{softplus}(x)=\log(1+e^x)$,
$\eta,\epsilon_s>0$, and $0<\epsilon_r<q_{\max}<1$ are fixed.
Subtracting the noise contribution can make $e_\ell$ negative;
the softplus transformation makes it positive and differentiable
for training.
The learned multiplier $e^{c_{\beta,\ell}}$ scales the positive
quantity in brackets, while the mapping bounds the output variance
$v_\ell q_\ell$ below the input variance $v_\ell$. This upper bound
ensures a positive variance after Gaussian message division.
Dividing this Gaussian output message by the input message gives
\begin{equation}
 \widetilde{\mathbf M}_{\ell+1}
 =\frac{\widehat{\mathbf X}_\ell-q_\ell\mathbf R_\ell}{1-q_\ell},
 \qquad
 \widetilde w_{\ell+1}=\frac{v_\ell q_\ell}{1-q_\ell}.
 \label{eq:twc-feedback-proposal}
\end{equation}
We damp the resulting message in its natural parameters
\cite{vehtari2020}:
\begin{equation}
 \begin{aligned}
 w_{\ell+1}^{-1}&=(1-\delta_\ell)w_\ell^{-1}
                    +\delta_\ell\widetilde w_{\ell+1}^{-1},\\
 \frac{\mathbf M_{\ell+1}}{w_{\ell+1}}
 &=(1-\delta_\ell)\frac{\mathbf M_\ell}{w_\ell}
       +\delta_\ell\frac{\widetilde{\mathbf M}_{\ell+1}}
                           {\widetilde w_{\ell+1}}.
 \end{aligned}
 \label{eq:twc-damping}
\end{equation}
Here $\delta_\ell=\min\{1,w_{\rm ref}/\widetilde w_{\ell+1}\}$
and $w_{\rm ref}=\rho P_H$, with
$P_H=\E_{q_{\rm tr}}\|h\|^2/D$ denoting the average channel power
per coefficient.
The damping reduces the weight of an outgoing message when its
variance exceeds $w_{\rm ref}$. Equations~\eqref{eq:twc-feedback-proposal}
and \eqref{eq:twc-damping} complete the update to layer $\ell+1$.
Through this feedback, both channel and activity estimates affect
the inputs to subsequent layers. The next section develops the
training objective for these coupled updates.
